\documentclass[10pt,a4paper]{amsart}
\usepackage[margin=19mm]{geometry}
\usepackage{amsmath,amssymb,mathtools}
\usepackage[scr=rsfs,cal=euler]{mathalfa}
\usepackage{xfrac}
\usepackage[unicode,hidelinks,bookmarks=false]{hyperref}
\newcommand{\ie}{i.e.\ }
\newcommand{\cf}{cf.\ }
\newcommand{\resp}{respectively\ }
\newcommand{\Subsection}{\subsection}
\providecommand{\nobreakdash}{\nobreak\hbox{-}}
\setlength{\emergencystretch}{2em}
\multlinegap0pt
\usepackage{bm}
\usepackage{bbm}
\usepackage{dsfont}
\usepackage{xspace}
\usepackage{cancel}
\usepackage{mathtools}
\usepackage{amsthm}
\makeatletter
\@ifundefined{lemma}{\newtheorem{lemma}{Lemma}}{}
\makeatother
\newcommand{\dd}{\mathrm{d}}
\title[Bayesian Best-Arm Identification with Abstention: A Polynomi]{Bayesian Best-Arm Identification with Abstention: A Polynomial-to-Exponential Phase Transition}
\author{Yuqi Huang, Yunlong Hou, Vincent Y. F. Tan}
\date{Source: arXiv:2606.29203v1 (CC BY 4.0)}
\begin{document}
\maketitle

\noindent\emph{Source and licence.} Bayesian Best-Arm Identification with Abstention: A Polynomial-to-Exponential Phase Transition. Version arXiv:2606.29203v1 (CC BY 4.0), distributed under the Creative Commons Attribution 4.0 International licence, as recorded in the arXiv licence metadata for this exact version and shown on the abstract page https://arxiv.org/abs/2606.29203v1. Full rights notice: licenses/arxiv-2606.29203-CC-BY-4.0.txt.\par\smallskip

\noindent\textit{Editorial scope.} Counted unit: the Lemma whose source label is \texttt{lem:one-arm-laplace-simple-upper} in the article (source0.tex lines 7352--7394, source label \texttt{lem:one-arm-laplace-simple-upper}), delivered together with the article's own complete original proof of it (source0.tex lines 7396--7696, one \texttt{proof} environment of 301 lines). No mathematics is written, repaired or re-derived: statement and proof are the article's own text line for line. Required context retained uncounted: eq:ckl-under-simple-upper (lines 6681--6687). Every definition, display and label of those lines is retained unchanged. Omitted from this package and not redistributed: 1--6680, 6688--7351, 7395--7395, 7697--8732. Nothing inside a retained excerpt is omitted or altered: every retained body is delivered exactly as the article's lines are written, so no omission receipt of any kind accompanies this package. Citations: the retained excerpts carry no citation command at all, so no bibliography entry of the article is transcribed and no reference is redistributed. Reference adaptation: none was needed and none was made; every reference inside the delivered unit resolves inside it, so no unresolved reference or citation is produced. Printed slips kept literal: no printed word, symbol, hypothesis or number of the counted statement or of its proof was altered. Layout and class: the article's own class is \texttt{article} with packages including enumerate, fontenc, amsmath, amssymb, natbib, xcolor, geometry, dsfont, thm-restate, pgfplots, smile\_new, multirow, rotating, esvect; the wrapper is the lane's pinned \texttt{amsart} editorial wrapper at 10pt on A4, which restates the theorem-like environments the retained text needs and adds the article's own macro definitions that the retained text needs (\texttt{\textbackslash dd}), with extra wrapper packages (bm, bbm, dsfont, xspace, cancel, mathtools). The delivered numbering is the unit's own. Assets: no retained span contains a figure environment, a graphics-inclusion command, a PDF-page inclusion, a table or a TikZ picture; no image file is copied into this package, the package declares no asset dependency and no image file is redistributed with it. Licence: version arXiv:2606.29203v1 of this article is distributed under the Creative Commons Attribution 4.0 International licence, as recorded in the arXiv licence metadata for this exact version and shown on the abstract page https://arxiv.org/abs/2606.29203v1. A full rights notice is delivered at licenses/arxiv-2606.29203-CC-BY-4.0.txt. Characters: every retained excerpt is pure ASCII, so the delivered engine, XeLaTeX with its default Latin Modern text font, reports no missing character.
\begin{equation}
\label{eq:ckl-under-simple-upper}
    D(\mu,z)
    \ge
    \widetilde{c}\,\{s(\mu)-s(z)\}^2,
    \qquad \mu,z\in\mathcal I .
\end{equation}

\begin{lemma}[One-arm Posterior Estimates for Ordering]
\label{lem:one-arm-laplace-simple-upper}
Fix an arm \(i\), and work on the event (that happens almost surely)
\[
    \frac1n\sum_{r=1}^nt(X_{i,r})
    \to
    A'(\theta(\mu_i)).
\]
Let \(\varPi_{i,n}\) be the one-arm posterior defined above.  Define
\[
    I_i(\eta):=D(\mu_i,s^{-1}(\eta)).
\]
Then the following hold.

First, for every compact interval \(K\Subset\mathcal H\),
\[
    \limsup_{n\to\infty}
    \frac1n\log\varPi_{i,n}(K)
    \le
    -\inf_{\eta\in K} I_i(\eta).
\]
Second, for every nonempty open interval \(G\subseteq\mathcal H\),
\[
    \liminf_{n\to\infty}
    \frac1n\log\varPi_{i,n}(G)
    \ge
    -\inf_{\eta\in G} I_i(\eta).
\]
Third, for every \(a,b\in\mathcal H\) with \(a<\eta_i<b\),
\[
    \limsup_{n\to\infty}
    \frac1n\log\varPi_{i,n}(\eta\le a)
    \le
    -I_i(a),
\]
and
\[
    \limsup_{n\to\infty}
    \frac1n\log\varPi_{i,n}(\eta\ge b)
    \le
    -I_i(b).
\]
\end{lemma}

\begin{proof}
Set
\[
    \eta_i=s(\mu_i),
    \qquad
    \theta_i:=\theta(\mu_i)=\beta(\eta_i),
    \qquad
    m_i:=A'(\theta_i),
\]
and write
\[
    \bar m_n:=\frac1n\sum_{r=1}^nt(X_{i,r}).
\]

For \(\eta\in\mathcal H\), write
\[
    p_\eta(w)
    :=
    \frac{\dd P_{s^{-1}(\eta)}}{\dd\lambda}(w)
    =
    \exp\{\beta(\eta)t(w)-A(\beta(\eta))\}h(w).
\]
Consider a realized full adaptive history $h_t=(a_1,x_1,\ldots,a_t,x_t)$,
For a candidate parameter vector
$\boldsymbol\eta=(\eta_1,\ldots,\eta_K)\in\mathcal H^K$,
the density of this realized history under the policy \(\pi\), with respect to
counting measure on actions and \(\lambda^{\otimes t}\) on rewards, is
\[
    L_t^\pi(\boldsymbol\eta;h_t)
    =
    \prod_{r=1}^t
    \pi_r(a_r\mid h_{r-1})
    \prod_{r=1}^t
    p_{\eta_{a_r}}(x_r).
\]
Here \(\pi_r(a_r\mid h_{r-1})\) is the sampling probability assigned by the
policy to the realized action \(a_r\) after the realized past \(h_{r-1}\).

Using the product prior in information coordinates, the posterior density of
\(\boldsymbol\eta\) is therefore proportional to
\[
    \prod_{\ell=1}^K
    \bar f_\ell(\eta_\ell)
    \prod_{r:a_r=\ell}
    p_{\eta_\ell}(x_r).
\]
Thus the posterior factorizes arm by arm.  In particular, the marginal posterior
of arm \(i\) depends only on the rewards observed from arm \(i\).  If these
rewards are denoted
$X_{i,1},\ldots,X_{i,n}$,
then, for every measurable set \(B\subseteq\mathcal H\),
\[
    \varPi_{i,n}(B)
    =
    \frac{
        \int_B
        \bar f_i(\eta)
        \prod_{r=1}^n p_\eta(X_{i,r})
        \,\dd\eta
    }{
        \int_{\mathcal H}
        \bar f_i(\eta)
        \prod_{r=1}^n p_\eta(X_{i,r})
        \,\dd\eta
    }.
\]
Substituting the exponential-family density gives
\[
    \prod_{r=1}^n p_\eta(X_{i,r})
    =
    \exp\left\{
        \beta(\eta)\sum_{r=1}^nt(X_{i,r})
        -
        nA(\beta(\eta))
    \right\}
    \prod_{r=1}^n h(X_{i,r}).
\]
\(\prod_{r=1}^n h(X_{i,r})\) does not depend on \(\eta\), so
it also cancels between numerator and denominator.
Bayes' rule gives
\[
    \varPi_{i,n}(B)
    =
    \frac{
        \int_B
        \bar f_i(\eta)
        \exp\{n[\beta(\eta)\bar m_n-A(\beta(\eta))]\}
        \,\dd\eta
    }{
        \int_{\mathcal H}
        \bar f_i(\eta)
        \exp\{n[\beta(\eta)\bar m_n-A(\beta(\eta))]\}
        \,\dd\eta
    }.
\]
Multiplying both numerator and denominator by
\[
    \exp\{-n[\beta(\eta_i)\bar m_n-A(\beta(\eta_i))]\}
\]
allows us to write
\[
    \varPi_{i,n}(B)
    =
    \frac{
        \int_B \bar f_i(\eta)e^{nL_n(\eta)}\,\dd\eta
    }{
        Z_n
    },
\]
where
\[
    L_n(\eta)
    :=
    (\beta(\eta)-\beta(\eta_i))\bar m_n
    -
    A(\beta(\eta))
    +
    A(\beta(\eta_i))
\]
and
\[
    Z_n
    :=
    \int_{\mathcal H}\bar f_i(\eta)e^{nL_n(\eta)}\,\dd\eta.
\]
Note \(L_n(\eta_i)=0\) and
\[
    L_n(\eta)\to (\beta(\eta)-\beta(\eta_i))A'(\beta(\eta_i)) - A(\beta(\eta))+A(\beta(\eta_i))
    = -D(\mu_i,s^{-1}(\eta))=-I_i(\eta)
\]
on compact subsets of \(\mathcal H\).



We first prove that
\begin{equation}
\label{eq:one-arm-denominator-rate-simple-upper}
    \frac1n\log Z_n\to0.
\end{equation}

For the lower bound on \(Z_n\), fix \(\tau>0\).  Since $\bar m_n\to m_i$, $L_n(\eta)\to -I_i(\eta)$ uniformly on compact neighborhoods of \(\eta_i\).  

Since \(I_i(\eta_i)=0\), choose a compact interval \(J_\tau\Subset\mathcal H\) containing \(\eta_i\) such that, for
all sufficiently large \(n\), $\inf_{\eta\in J_\tau}L_n(\eta)\ge -\tau$.

By continuity and strict positivity of \(\bar f_i\),
\[
    \int_{J_\tau}\bar f_i(\eta)\,\dd\eta>0.
\]
Therefore
\[
    Z_n
    \ge
    \int_{J_\tau}\bar f_i(\eta)e^{nL_n(\eta)}\,\dd\eta
    \ge
    e^{-n\tau}\int_{J_\tau}\bar f_i(\eta)\,\dd\eta.
\]
Thus
\[
    \liminf_{n\to\infty}\frac1n\log Z_n\ge -\tau.
\]
Letting \(\tau\downarrow0\) gives
\[
    \liminf_{n\to\infty}\frac1n\log Z_n\ge0.
\]

For the upper bound on \(Z_n\), for all large \(n\) we have
\(\bar m_n\in A'(\Theta)\).  Define the pseudo-MLE
\[
    \theta_n:=(A')^{-1}(\bar m_n),
    \qquad
    \eta_n:=\beta^{-1}(\theta_n).
\]
Then $\theta_n\to\theta_i$ and $\eta_n\to \eta_i$.
Moreover, using \(\bar m_n=A'(\theta_n)\),
\[
\begin{aligned}
    L_n(\eta_n)-L_n(\eta)
    &=
    (\theta_n-\beta(\eta))A'(\theta_n)
    -
    A(\theta_n)
    +
    A(\beta(\eta))                                                \\
    &=
    D(s^{-1}(\eta_n),s^{-1}(\eta)).
\end{aligned}
\]
Therefore, by assumption \eqref{eq:ckl-under-simple-upper},
\[
    L_n(\eta)
    \le
    L_n(\eta_n)-\widetilde{c}\,(\eta-\eta_n)^2.
\]
Since
\[
    L_n(\eta_n)
    =
    D(s^{-1}(\eta_n),\mu_i)
    \to0,
\]
boundedness of \(\bar f_i\) gives
\[
\begin{aligned}
    Z_n
    &\le
    \|\bar f_i\|_\infty
    e^{nL_n(\eta_n)}
    \int_{\mathcal H}e^{-n\widetilde{c} (\eta-\eta_n)^2}\,\dd\eta        \\
    &\le
    \|\bar f_i\|_\infty
    e^{nL_n(\eta_n)}
    \sqrt{\frac{\pi}{n\widetilde{c}}}.
\end{aligned}
\]
Hence
\[
    \limsup_{n\to\infty}\frac1n\log Z_n\le0.
\]
This proves \eqref{eq:one-arm-denominator-rate-simple-upper}.

Now fix a compact interval \(K\Subset\mathcal H\).  Since \(\bar m_n\to m_i\),
\[
    L_n(\eta)\to -D(\mu_i,s^{-1}(\eta))=-I_i(\eta)
\]
uniformly on \(K\).  Hence
\[
\begin{aligned}
    \limsup_{n\to\infty}
    \frac1n\log
    \int_K \bar f_i(\eta)e^{nL_n(\eta)}\,\dd\eta
    &\le
    -\inf_{\eta\in K}I_i(\eta).
\end{aligned}
\]
Combining this numerator bound with
\eqref{eq:one-arm-denominator-rate-simple-upper} gives the compact-interval
upper bound.

Next let \(G\subseteq\mathcal H\) be a nonempty open interval.  Choose
\(\eta_0\in G\) such that
\[
    I_i(\eta_0)\le \inf_{\eta\in G}I_i(\eta)+\tau.
\]
Choose a compact interval \(J\Subset G\) containing \(\eta_0\) and small enough
that
\[
    \sup_{\eta\in J}I_i(\eta)\le I_i(\eta_0)+\tau.
\]
By uniform convergence of \(L_n\) on \(J\) and strict positivity of
\(\bar f_i\) on \(J\),
\[
\begin{aligned}
    \liminf_{n\to\infty}
    \frac1n\log
    \int_G \bar f_i(\eta)e^{nL_n(\eta)}\,\dd\eta
    &\ge
    \liminf_{n\to\infty}
    \frac1n\log
    \int_J \bar f_i(\eta)e^{nL_n(\eta)}\,\dd\eta                  \\
    &\ge
    -\sup_{\eta\in J}I_i(\eta)
    \ge
    -\inf_{\eta\in G}I_i(\eta)-2\tau.
\end{aligned}
\]
Using again the denominator rate
\eqref{eq:one-arm-denominator-rate-simple-upper}, and then letting
\(\tau\downarrow0\), gives the open-interval lower bound.

It remains to prove the one-sided estimates.  Consider \(a<\eta_i\).  Since
\(\eta_n\to \eta_i\), eventually \(a<\eta_n\).  The function $\theta\mapsto \theta\bar m_n-A(\theta)$
is strictly concave and maximized at \(\theta_n\).  Since \(\beta\) is strictly
increasing, \(L_n(\eta)\) is increasing on \(\{\eta<\eta_n\}\).  Therefore, for all large
\(n\),
\[
    \sup_{\eta\le a}L_n(\eta)=L_n(a).
\]
Thus
\[
    \varPi_{i,n}(\eta\le a)
    \le
    \frac{e^{nL_n(a)}}{Z_n},
\]
because \(\int_{\{\eta\le a\}}\bar f_i(\eta)\,\dd\eta\le1\).
Since $L_n(a)\to -I_i(a)$
and \(n^{-1}\log Z_n\to0\), we obtain
\[
    \limsup_{n\to\infty}
    \frac1n\log\varPi_{i,n}(\eta\le a)
    \le
    -I_i(a).
\]
The proof for \(b>\eta_i\) is identical: eventually \(b>\eta_n\), \(L_n\) is
decreasing on \(\{\eta>\eta_n\}\), and hence
\[
    \sup_{\eta\ge b}L_n(\eta)=L_n(b).
\]
This gives the upper-tail estimate.

\end{proof}

\end{document}
